\scanpage{200}\mspage{149}
Comme $X_{0y}\longrightarrow X_{1y}$ est un iso donc est dans $W$, s'ensuit que $e\to\unsure{\mathscr X|\mathcal X}_y$ l'est aussi (prop. 13 p. 125), moyennant $W(7\mathrm{bis})$, donc $\mathscr X_y$ est $W$-asphérique, OK.

% slanted
\margin{On n'a besoin ici que de la forme directe de $W(7\mathrm{bis})$, i.e (notations de \uline{XVI}) de $\mathrm{Loc}(6'\mathrm a)$.}

\uline{Corollaire 3} \inserted{(Moyennant $W(7\mathrm{bis})$)}\margin{inutile} Soit $X$ dans $\mathrm{Cat}$, $U_1,U_2$ deux ouverts de $X$ tels que $X=U_1\cup U_2$, $U_0=U_1\cap U_2$ l'ouvert intersection. Alors le carré d'imm. ouvertes

\diagram{\begin{tikzcd}[ampersand replacement=\&,row sep=1.5em,column sep=2em]
U_0\arrow[r,hook]\arrow[d,hook] \& U_1\arrow[d,hook]\\
U_2\arrow[r,hook] \& X
\end{tikzcd}\(\,.\)}

est $W$-cocartésien. (\unsure{Énoncé} dual pour deux fermés $T_1,T_2$ tels que $T_1\cup T_2=\unsure{x|X}$, avec $T_0=T_1\cap T_2$\ldots)

On applique le cor. 2. On sait que le carré est cocartésien dans $\mathrm{Cat}$${}^{*}$, de plus $U_1\hookrightarrow X$ est bien $W$-lisse, \deleted{\unsure{étant}} une imm. ouverte.

% slanted
\margin{On n'a pas même besoin de le savoir, cf annotation du cor. 2}

\uline{Corollaire 4} \inserted{(Moyennant $W(7\mathrm{bis})^{*}$)} Soit $U\overset{i}{\hookrightarrow}X$ une imm. ouverte telle que $T=X-U$ soit isom. à la catégorie ponctuelle $e$. Pour que $i\in W$, il f. et s. que $i$ soit $W$-asphérique, i.e. que $U_{/t}$ \deleted{soit} (où $t$ est l'unique objet de $T$) soit $W$-asphérique.

% slanted
\margin{Ici on \unsure{a} besoin de $W(7\mathrm{bis})$ forme forte i.e. $i\circ r\in W\unsure{\Longleftarrow|\Longleftrightarrow}i\in W$, cf annotation Th. 4.}

Il suffit bien sûr de prouver que \deleted{\unsure{$i$}}

\[
i\in W\Longrightarrow U_{/t}\ \text{$W$-asphérique}.
\]
% archive sheet number: 200
